\subsection{Dynamic $\alpha$-Scheduling}
\label{sec:alpha-scheduling}

A fixed $\alpha$ ignores that the value of exploration should shrink as the
campaign progresses, echoing the broader practice of decaying exploration
coefficients over time in the bandit and reinforcement-learning literature
\citep{auer2002finite}. We schedule $\alpha(t)$ against two natural progress
variables: $\mu = N(t)/N$, the fraction of the sending budget already used,
and $\tilde o = o(t)/m$, the fraction of recipients who have already
opened, with $o(t) = x(t)^\top\mathbf{1}$.

\paragraph{Send-volume-based (function of $\mu$).}
A \textbf{linear} schedule from $l$ to $r$, in the spirit of the linearly
annealed $\varepsilon$-greedy exploration rate used to train DQN agents
\citep{mnih2015human},
\begin{equation}
\label{eq:alpha-linear}
\alpha(\mu) = l(1-\mu) + r\mu,
\end{equation}
and a \textbf{geometric} (log-linear) schedule, analogous to the geometric
cooling schedules of simulated annealing \citep{kirkpatrick1983optimization},
\begin{equation}
\label{eq:alpha-geometric}
\alpha(\mu) = l^{1-\mu}\, r^{\mu}.
\end{equation}

\paragraph{Open-rate-based (function of $\tilde o$).}
Since $f_j(t)$ is itself a point estimate that becomes more reliable as
more opens are observed, we also consider a \textbf{confidence} schedule,
hyperbolic in $\tilde o$ --- paralleling how the confidence radius of
UCB-style bandit policies shrinks as observations accumulate
\citep{auer2002finite}:
\begin{equation}
\label{eq:alpha-confidence}
\alpha(\tilde o) = \frac{\kappa}{\tilde o + \kappa},
\end{equation}
with dampening factor $\kappa$ (so $\alpha(\kappa)=\tfrac12$). All three
reduce to the constant baseline at their degenerate limits ($l=r$, or
$\kappa\to\infty$). A schedule jointly conditioned on both $\mu$ and
$\tilde o$ is conceptually natural but is left to future work
(Section~\ref{sec:future-work}) to avoid an unconstrained hyperparameter
count.
